\chapter{Local coefficients and ordered collisions}
\label{ch:vol2-coefficients}

When insertion points collide, an operation can acquire a pole or a
term supported on the collision. Both depend on the normal coordinate,
but they obey different multiplication rules. A pole permits division
by that coordinate. A supported term can be annihilated by it.
Their relation is an attachment between coefficient modules.
To place the comparison homotopies of Chapter~\ref{ch:vol2-boundary-actions}
on ordered configurations, we first construct these modules and their
derived maps.

\section{Principal parts and their residue currents}
\label{sec:vol2-principal-parts}

Work over $\C$, with one algebraic normal coordinate $z$. Put
\[
 A=\C[z],\qquad L=\C[z,z^{-1}],\qquad E=L/A.
\]
The quotient $E$ is the module of principal parts at the origin.
It has basis $t_n=z^{-n}\bmod A$, $n\geq1$, over $\C$.
Every element is a finite linear combination of these basis vectors.

Multiplication and differentiation give
\begin{equation}\label{eq:vol2-normal-jets}
 zt_1=0,\qquad zt_n=t_{n-1}\quad(n>1),\qquad
 \partial_z t_n=-nt_{n+1}.
\end{equation}
Thus $E$ retains derivatives of every normal order, even though each
of its elements is killed by a power of $z$.

Let $D=A\langle\partial_z\rangle$ be the algebra with relation
$\partial_z z-z\partial_z=1$. A left $D$-module is an $A$-module
with an operator $\nabla$ satisfying
$\nabla(am)=a'm+a\nabla m$. Such an operator is its connection
in the coordinate $z$. The usual derivative makes
\begin{equation}\label{eq:vol2-localization-sequence}
 0\longrightarrow A\longrightarrow L\longrightarrow E\longrightarrow0
\end{equation}
an exact sequence of left $D$-modules.

\begin{proposition}\label{prop:vol2-normal-module}
The map $D/Dz\to E$ sending the class of $1$ to $t_1$ is an
isomorphism of left $D$-modules. There is no unital associative algebra
extending the multiplication on $L$ and containing a nonzero element
$T$ with $zT=0$.
\end{proposition}

\begin{proof}
Using $z\partial_z=\partial_z z-1$, every element of $D/Dz$
has a representative that is a polynomial in $\partial_z$.
Its monomials act on $t_1$ by
$\partial_z^r t_1=(-1)^r r!t_{r+1}$. These images form a basis of $E$,
so the map is both injective and surjective.
In the proposed algebra, associativity and the inverse of $z$ give
$T=z^{-1}(zT)=0$, a contradiction.
\end{proof}

The relation with an analytic contact term fixes the normalization.
Write $z=x+iy$, orient the plane by $dx\wedge dy$, and set
$\partial_z=\tfrac12(\partial_x-i\partial_y)$ and
$\partial_{\bar z}=\tfrac12(\partial_x+i\partial_y)$.

A distribution acts on compactly supported smooth test functions.
Its derivatives are defined by integration by parts. Let $\delta_0$
be the distribution $\langle\delta_0,\phi\rangle=\phi(0)$.
The locally integrable function $1/z$ defines a distribution $P_1$.
For $n\geq1$, define its higher principal values by
\[
 P_n=\frac{(-1)^{n-1}}{(n-1)!}\partial_z^{n-1}P_1.
\]
They agree with $z^{-n}$ off the origin.

\begin{proposition}\label{prop:vol2-residue-current}
With these conventions,
\begin{equation}\label{eq:vol2-residue-normalization}
 \bar\partial P_n
 =\frac{\pi(-1)^{n-1}}{(n-1)!}
       \partial_z^{n-1}\delta_0\,d\bar z.
\end{equation}
The assignment of the right-hand side to $t_n$ identifies $E$ with
the left $D$-module spanned by these residue currents. The coefficient
module is placed in cochain degree zero. Its currents have Dolbeault
degree one when used in a Dolbeault complex.
\end{proposition}

\begin{proof}
For a test function $\phi$, Stokes' formula on the complement of the
disc of radius $\varepsilon$ gives
\[
 -\lim_{\varepsilon\to0}\int_{|z|>\varepsilon}
       \frac{\partial_{\bar z}\phi}{z}\,dx\,dy
 =\lim_{\varepsilon\to0}\frac1{2i}
       \int_{|z|=\varepsilon}\frac{\phi(z)}z\,dz
 =\pi\phi(0).
\]
The circle on the middle expression is counterclockwise. The outer
boundary contributes zero because $\phi$ is compactly supported.
This proves $\bar\partial P_1=\pi\delta_0d\bar z$.
Distributional derivatives commute, giving
\eqref{eq:vol2-residue-normalization} for every $n$.

The identity $z\partial_z^r\delta_0=-r\partial_z^{r-1}\delta_0$
follows by differentiating $z\delta_0=0$. It gives precisely the
multiplication rule in \eqref{eq:vol2-normal-jets}. Differentiation
gives its connection rule. The derivatives of $\delta_0$
are linearly independent: test them against $z^r$ times a smooth
function equal to one near the origin. Thus the assignment is an
isomorphism onto the displayed span.
\end{proof}

The residue is therefore a map between specified coefficient objects.
It supplies no multiplication of a principal value by an arbitrary
supported distribution. To compare operations using these objects,
we need resolutions that retain their connections.

\section{Resolutions and derived coefficient maps}
\label{sec:vol2-coefficient-resolutions}

Complexes have cohomological grading, with $V[1]^j=V^{j+1}$.
A complex of differential modules has a connection commuting with
its differential.

Define a free $A$-complex $Q_L$ by
\begin{equation}\label{eq:vol2-coefficient-telescope}
 Q_L^{-1}=\bigoplus_{n\geq0}A\alpha_n,\qquad
 Q_L^0=\bigoplus_{n\geq0}Ae_n,\qquad
 d\alpha_n=e_n-ze_{n+1}.
\end{equation}
Its augmentation sends $e_n$ to $z^{-n}$. Give it the connection
\begin{equation}\label{eq:vol2-telescope-connection}
 \nabla e_n=-ne_{n+1},\qquad
 \nabla\alpha_n=-(n+1)\alpha_{n+1},
\end{equation}
extended by the Leibniz rule. Put $Q_A=A$ in degree zero, and let
$r:Q_A\to Q_L$ send $1$ to $e_0$. Set $Q_E=\operatorname{Cone}(r)$,
with the convention
\begin{equation}\label{eq:vol2-coefficient-cone}
 \operatorname{Cone}(r)^j=Q_L^j\oplus Q_A^{j+1},\qquad
 d(\zeta,\xi)=(d_{Q_L}\zeta+r\xi,-d_{Q_A}\xi).
\end{equation}

\begin{lemma}\label{lem:vol2-coefficient-telescope}
The augmentations $Q_L\to L$ and $Q_E\to E$ are
quasi-isomorphisms of differential modules. All terms of $Q_A,Q_L,Q_E$
are free over $A$.
\end{lemma}

\begin{proof}
If a finite sum $\sum a_n(e_n-ze_{n+1})$ vanishes, its coefficient
of $e_0$ gives $a_0=0$. The successive coefficients give $a_n=0$
for every $n$. Thus the differential in
\eqref{eq:vol2-coefficient-telescope} is injective.
Its cokernel imposes $e_n=ze_{n+1}$ and is the direct limit
of $A\xrightarrow{z}A\xrightarrow{z}\cdots$, identified with $L$
by $e_n\mapsto z^{-n}$.

The connection commutes with the differential because
\[
 \nabla(e_n-ze_{n+1})
 =-(n+1)e_{n+1}+(n+1)ze_{n+2}
 =d\bigl(-(n+1)\alpha_{n+1}\bigr).
\]
It also commutes with both augmentations and with $r$.
The cone of the injection $A\to L$ has cohomology $E$ in degree
zero and no other cohomology. The cone in
\eqref{eq:vol2-coefficient-cone} consequently resolves $E$.
Its terms are direct sums of the indicated free modules.
\end{proof}

These resolutions are free over $A$. To obtain derived maps of
differential modules, we also resolve the derivative action.
For any one of the three complexes $Q$, use the two Spencer columns
\begin{equation}\label{eq:vol2-spencer}
 D\otimes_A Q\xrightarrow{s}D\otimes_A Q,\qquad
 s(P\otimes q)=P\partial_z\otimes q-P\otimes\nabla q,
\end{equation}
in degrees $-1,0$. In column degree $p$, the total differential is
$s+(-1)^p d_Q$. Denote this total complex by $P_Q$.
The Leibniz rule makes $s$ well defined under $A$-balancing.

\begin{proposition}\label{prop:vol2-spencer-maps}
The complexes $P_{Q_A},P_{Q_L},P_{Q_E}$ are bounded free
$D$-resolutions of $A,L,E$. Their complexes of $D$-linear maps,
with ordinary graded composition, form a differential graded category
computing all derived maps among these three objects.

For any complex $N$ of left $D$-modules, put
$C=\Hom_A(Q,N)$, using product totalization, and define
$\mathcal Df=\nabla_Nf-f\nabla_Q$. Then $\RHom_D(Q,N)$ is represented
by the complex with degree-$j$ term $C^j\oplus C^{j-1}$ and differential
\begin{equation}\label{eq:vol2-connection-hom}
 d(a,b)=(d_Ca,\mathcal Da-d_Cb),\qquad
 d_Cf=d_Nf-(-1)^{|f|}fd_Q.
\end{equation}
\end{proposition}

\begin{proof}
For a differential module $M$, augment the second column of
\eqref{eq:vol2-spencer} by $P\otimes m\mapsto Pm$.
Write $D$ in its right $A$-basis $1,\partial_z,\partial_z^2,\ldots$.
The highest derivative of a nonzero finite sum acquires a nonzero
term of one higher order under $s$, so $s$ is injective.
Its relations successively replace $\partial_z^r\otimes m$ by
$1\otimes\nabla^r m$. Every class in its cokernel therefore has a
unique representative $1\otimes m$, proving exactness.

Apply this resolution in each degree of $Q$. The double complex
is bounded in both degrees, so the total augmentation $P_Q\to Q$
is a quasi-isomorphism. Its terms are free over $D$ because the
terms of $Q$ are free over $A$. A bounded complex of projective
modules computes derived Hom, including when its free terms have
countably many generators.

The adjunction
$\Hom_D(D\otimes_AQ,N)=\Hom_A(Q,N)$ identifies the map induced by
$s$ with $\mathcal D$. With the sign of the second summand chosen as
in \eqref{eq:vol2-connection-hom}, the total Hom differential is
the one displayed there. For every $A$-linear $f$, the map
$\mathcal Df$ is again $A$-linear. The operator $\mathcal D$ is
$\C$-linear, obeys $\mathcal D(af)=a'f+a\mathcal Df$ for $a\in A$,
and commutes with $d_C$. Hence the displayed differential squares
to zero.

The direct sums in $Q$ give products in its Hom complexes.
Using the free $D$-resolutions on both sides supplies associative
composition and the graded Leibniz identity.
\end{proof}

The extra column in \eqref{eq:vol2-connection-hom} records compatibility
with normal differentiation. The underlying $A$-linear Hom complex
alone cannot express that compatibility. A second distinction appears
when two supported coefficients meet: their tensor product must also
be derived.

\section{Supported tensors and redundant equations}
\label{sec:vol2-supported-tensors}

Write $D(A)$ for the derived category of $A$-modules. The symbol
$\otimes_A^{\mathbb L}$ means the tensor computed after a flat
resolution. All tensors in the next proposition are over $A$, rather
than over $D$. Let $k_0=A/(z)$ be the coherent closed fibre.
It has no left $D$-action extending its $A$-action: since $z$ acts
by zero, the commutator $[\partial_z,z]$ could not act by one.
The differential module supported at the origin is instead $E$,
as in Proposition~\ref{prop:vol2-normal-module}.

\begin{proposition}\label{prop:vol2-supported-self-tensor}
In $D(A)$ one has
\begin{equation}\label{eq:vol2-supported-self-tensor}
 E\otimes_AE=0,\qquad E\otimes_A^{\mathbb L}E\simeq E[1],
 \qquad
 k_0\otimes_A^{\mathbb L}k_0\simeq k_0\otimes\Lambda(\epsilon),
 \quad |\epsilon|=-1.
\end{equation}
The last expression has zero differential.
\end{proposition}

\begin{proof}
Every element of $E$ is $z$-power torsion, whereas multiplication by
$z$ on $E$ is surjective. For a simple tensor $e\otimes e'$,
choose $m$ with $z^me=0$ and write $e'=z^mf$. Then
$e\otimes e'=z^me\otimes f=0$.

The flat resolution $[A\to L]$ of $E$ is in degrees $-1,0$.
Since $L\otimes_AE=0$, its tensor with $E$ is $[E\to0]$,
which is $E[1]$.

For the closed fibre, tensor the free resolution
$[A\xrightarrow{z}A]$ with $k_0$. Its differential becomes zero.
The Koszul multiplication gives the exterior generator $\epsilon$.
\end{proof}

Several collision equations introduce a further choice. For a
commutative ring $B$ and an element $f\in B$, let
$B_f=B[f^{-1}]$ and define
\[
 C_f=[B\longrightarrow B_f]
\]
in degrees $0,1$. For a finite list $\mathbf f=(f_1,\ldots,f_s)$,
its \emph{stable support complex} is
\begin{equation}\label{eq:vol2-stable-support}
 C_{\mathbf f}=C_{f_1}\otimes_B\cdots\otimes_B C_{f_s},
\end{equation}
with the usual tensor differential. Its terms are flat. The
projection $C_f\to B$ is the identity in degree zero and zero
in degree one.

\begin{lemma}\label{lem:vol2-redundant-support}
If $g$ lies in the radical of the ideal $(f_1,\ldots,f_s)$, then
the projection
\[
 C_{\mathbf f}\otimes_B C_g\longrightarrow C_{\mathbf f}
\]
is a quasi-isomorphism.
\end{lemma}

\begin{proof}
The complex $C_f$ is the filtered direct limit of
$[B\xrightarrow{f^n}B]$, $n\geq1$. The transition is the identity
in degree zero and multiplication by $f$ in degree one.
Consequently $C_{\mathbf f}$ is the direct limit of
$\bigotimes_i[B\xrightarrow{f_i^n}B]$.
On the $i$th two-term factor, the identity from degree one to
degree zero is a homotopy for multiplication by $f_i^n$.
Its signed tensor with the identities on the other factors gives
the same homotopy on the full tensor product.
Filtered direct limits of modules are exact, so each cohomology
class of $C_{\mathbf f}$ is annihilated by some common powers
$f_1^n,\ldots,f_s^n$.

Choose $N$ with $g^N\in(f_1,\ldots,f_s)$. The ideal
$(f_1,\ldots,f_s)^{s(n-1)+1}$ is contained in
$(f_1^n,\ldots,f_s^n)$, since each generating monomial contains
one $f_i$ at least $n$ times. A power of $g$ therefore annihilates
the given class. It follows that $C_{\mathbf f}\otimes_B B_g$
is acyclic, since localization is exact. Tensoring the triangle
$C_g\to B\to B_g$ with the flat complex $C_{\mathbf f}$ proves
the assertion.
\end{proof}

The finite Koszul complex retains different information. Take
$B=\C[u,v]$ and a third equation $w=u+v$. The finite complex
$K(u,v,w)$ is the exterior algebra over $B$ on degree-$-1$
generators $e_u,e_v,e_w$, with differential
\[
 de_u=u,\qquad de_v=v,\qquad de_w=u+v.
\]
The change of basis $\eta=e_w-e_u-e_v$ gives
\begin{equation}\label{eq:vol2-triangle-excess}
 K(u,v,u+v)=K(u,v)\otimes\Lambda(\eta),\qquad
 d\eta=0,
\end{equation}
and hence $H^*K(u,v,u+v)=\C\otimes\Lambda(\eta)$.
Indeed, $u,v$ is a regular sequence: multiplication by $u$ is
injective in $B$, and multiplication by $v$ is injective in $B/(u)$.
The two successive two-term resolutions prove that $K(u,v)$
resolves $B/(u,v)=\C$.

By contrast, $C_{u,v,u+v}\simeq C_{u,v}$ by
Lemma~\ref{lem:vol2-redundant-support}. The latter has only
degree-two cohomology,
\begin{equation}\label{eq:vol2-two-normal-parts}
 H^2C_{u,v}=
 E_2=\frac{\C[u^{\pm1},v^{\pm1}]}
              {\C[u^{\pm1},v]+\C[u,v^{\pm1}]},
\end{equation}
with basis $u^{-m}v^{-n}$, $m,n\geq1$.
To see the concentration, write $C_{u,v}$ as the tensor over $\C$
of the one-variable complexes for $u$ and $v$. Each has its
principal-parts cohomology in degree one.

The vanishing of $(E_2)_{u+v}$ also has a direct calculation:
\[
 \frac1{u^mv^n}=
 \frac1{u+v}\left(\frac1{u^{m-1}v^n}
                       +\frac1{u^mv^{n-1}}\right).
\]
Repeated application decreases $m+n$ until every term is a power
of $(u+v)^{-1}$ times a fraction regular in one normal coordinate.
That fraction represents zero in $E_2$, so the term vanishes in
$(E_2)_{u+v}$.
The odd class $\eta$ in \eqref{eq:vol2-triangle-excess} belongs
to the finite closed-intersection complex. It does not supply an
extra class in the stable support complex.

\section{Differential modules on configuration spaces}
\label{sec:vol2-geometric-coefficients}

The coordinate calculations extend over a smooth separated complex
algebraic curve $X$. For a smooth scheme $Y$, let $\Dcal_Y$ be its
sheaf of differential operators. Let $\operatorname{DMod}(Y)$
denote the stable enhancement of the unbounded derived category
of right $\Dcal_Y$-modules with quasi-coherent underlying
$\Ocal_Y$-modules. The enhancement retains mapping complexes
and coherent homotopies. We use the enhanced differential-module
formalism of Francis--Gaitsgory \cite{FG12}, \S1.4, including
extraordinary inverse image, direct image and external product,
with their composition coherences.

Right modules give a coordinate-independent form of our local
left modules. If $M$ is a left differential module, the associated
right module is $\omega_Y\otimes_{\Ocal_Y}M$, where $\omega_Y$
is the line of top forms. In one coordinate its action satisfies
\[
 (dz\otimes m)\partial_z=-dz\otimes\nabla m.
\]
For a nonconstant form, the derivative also acts on its coefficient.
This is the transposition that fixes functions, negates coordinate
derivatives and reverses products. It preserves the normal derivatives
in $E$.

Here are the functors and shifts used below. For a map $f:Z\to Y$
of smooth schemes, with dimensions $d_Z,d_Y$, the transfer modules are
\[
 \Dcal_{Z\to Y}
 =\Ocal_Z\otimes_{f^{-1}\Ocal_Y}f^{-1}\Dcal_Y,
 \qquad
 \Dcal_{Y\leftarrow Z}
 =\omega_Z\otimes_{\Ocal_Z}\Dcal_{Z\to Y}
       \otimes_{f^{-1}\Ocal_Y}f^{-1}\omega_Y^{-1}.
\]
The first has a left $\Dcal_Z$-action and a right
$f^{-1}\Dcal_Y$-action. Transposition gives the second a left
$f^{-1}\Dcal_Y$-action and a right $\Dcal_Z$-action.

For right-module complexes $M$ on $Z$ and $N$ on $Y$, set
\begin{align}
 f_+M&=Rf_*^{\Ocal}
       (M\otimes_{\Dcal_Z}^{\mathbb L}\Dcal_{Z\to Y}),
       \label{eq:vol2-dmodule-direct}\\
 f^!N&=(f^{-1}N\otimes_{f^{-1}\Dcal_Y}^{\mathbb L}
                  \Dcal_{Y\leftarrow Z})[d_Z-d_Y].
       \label{eq:vol2-dmodule-inverse}
\end{align}
Here $f^{-1}$ is inverse image of sheaves, and $Rf_*^{\Ocal}$
is derived sheaf direct image. External product, denoted $\boxtimes$,
places two differential modules on independent factors.

For the structure map $p_Y:Y\to\Spec\C$, our normalization is
\begin{equation}\label{eq:vol2-dualizing-shift}
 J_Y=p_Y^!\C=\omega_Y[d_Y],\qquad
 J_Y\boxtimes J_Z=J_{Y\times Z}.
\end{equation}

For an open immersion $j$, write $j_*$ for $j_+$. Its functor
$j^!$ is restriction. For a closed immersion $i$, write $i_*=i_+$.
The functor $i_*$ retains all normal derivatives along the closed
subspace.

Differential modules on a possibly singular reduced
closed subspace are defined by the corresponding supported modules
in smooth embeddings, using Kashiwara's equivalence. Thus the
open--closed recollement identities apply to unions of diagonals:
\begin{equation}\label{eq:vol2-open-recollement}
 i_*i^!N\longrightarrow N\longrightarrow j_*j^!N,
 \qquad i^!j_*=0,\qquad j^!j_*=\id.
\end{equation}
These identities, and the coherent composition laws of the functors,
are part of the specified coefficient formalism.

The cut product will first restrict to a disjoint locus and then
apply $j_*$. Its compatibility with a collision rests on the
following consequence of recollement.

\begin{lemma}\label{lem:vol2-open-basechange}
For a Cartesian square with $j$ an open immersion,
\[
 \begin{tikzcd}
 U'\arrow[r,"f_U"]\arrow[d,"j'"']&U\arrow[d,"j"]\\
 Y'\arrow[r,"f"']&Y,
 \end{tikzcd}
 \qquad
 f^!j_*M\simeq j'_*f_U^!M.
\]
These equivalences are coherent under composition of $f$.
\end{lemma}

\begin{proof}
Let $i:Z\hookrightarrow Y$ and $i':Z'\hookrightarrow Y'$
be the reduced closed complements. The map $fi'$ factors as $ih$.
Transitivity and \eqref{eq:vol2-open-recollement} give
\[
 i'^!f^!j_*M=h^!i^!j_*M=0.
\]
The recollement triangle on $Y'$ therefore identifies $f^!j_*M$
with $j'_*j'^!f^!j_*M$. On the open set $U'$, transitivity
identifies its restriction with $f_U^!M$.

For a composite $f$, all such comparisons lie in the fully faithful
image of $j'_*$. Their restrictions are the given coherent
composition maps for extraordinary inverse image. Full faithfulness
lifts those maps and all their higher homotopies uniquely.
\end{proof}

\section{Ordered words and consecutive cuts}
\label{sec:vol2-word-coefficients}

An ordered insertion retains its place even when its point coincides
with another point. The word $(x,y,x)$, with $x\ne y$, is an elementary
example. Its two consecutive cuts are $(x)\mid(y,x)$ and
$(x,y)\mid(x)$. Neither separates disjoint sets of points.

The intervening operation cannot be deleted. For the matrix units
$E_{ij}$ on $\C^2$, which send the $j$th basis vector to the $i$th,
\[
 E_{12}E_{21}E_{12}=E_{12},\qquad E_{12}^2=0.
\]
We accordingly retain words and allow adjacent positions to collapse
without identifying a word with its unordered set of values.

Let $[n]=\{1,\ldots,n\}$ with its usual order, and let
$\operatorname{OrdSurj}_+$ be the category of these nonempty ordered
sets and monotone surjections. A surjection
$\alpha:[n]\twoheadrightarrow[m]$ defines
\[
 d_\alpha:X^m\longrightarrow X^n,\qquad
 d_\alpha(x_1,\ldots,x_m)_i=x_{\alpha(i)}.
\]
Write $D(\C)$ for complexes of complex vector spaces up to
quasi-isomorphism, with their stable enhancement. The category of
\emph{word coefficients} is
\begin{equation}\label{eq:vol2-word-category}
 \mathcal W_X=D(\C)\times
   \varprojlim_{\operatorname{OrdSurj}_+}\operatorname{DMod}(X^n),
\end{equation}
where the transition from $n$ to $m$ is $d_\alpha^!$.
An object $F$ consists of an empty coefficient $F_0$, coefficients
$F_n$ on $X^n$, equivalences $d_\alpha^!F_n\simeq F_m$, and
their coherent compatibility under composition. The empty
coefficient has no collision arrow to a nonempty one.

For $0<p<n$, define the cut open and its inclusion by
\begin{equation}\label{eq:vol2-cut-open}
 U_{p,n}=\{(x_1,\ldots,x_n):x_i\ne x_j
                         \text{ whenever }i\leq p<j\},
 \qquad j_{p,n}:U_{p,n}\hookrightarrow X^n.
\end{equation}
For $p=0,n$, put $U_{p,n}=X^n$ and let $j_{p,n}$ be the identity.
Define
\begin{equation}\label{eq:vol2-cut-product}
 (F\star_{\mathrm{cut}}G)_n
 =\bigoplus_{p=0}^n j_{p,n,*}j_{p,n}^!
                  (F_p\boxtimes G_{n-p}).
\end{equation}
An empty factor means scalar tensoring. In particular,
$(F\star_{\mathrm{cut}}G)_0=F_0\otimes_\C G_0$.

\begin{theorem}\label{thm:vol2-cut-product}
Formula~\eqref{eq:vol2-cut-product} defines an associative monoidal
product on $\mathcal W_X$. Its unit $K$ has $K_0=\C$ and
$K_n=0$ for $n>0$. Its associativity maps are the identifications
of the iterated products with the sums over consecutive partitions
whose blocks have pairwise disjoint values in $X$.
\end{theorem}

\begin{proof}
First pull a summand of \eqref{eq:vol2-cut-product} back along
$d_\alpha$. If its cut lies inside a fibre of $\alpha$, two
positions on opposite sides of the cut are forced to have the same
value. Its inverse-image open is empty, and
Lemma~\ref{lem:vol2-open-basechange} makes the summand zero.
Otherwise the cut lies between fibres, including either end, hence
determines a unique cut $q$ of $[m]$. Its inverse-image open is
$U_{q,m}$. External-product compatibility, the word equivalences
of $F,G$, and the same lemma identify the pulled-back summand
with the $q$th summand of $(F\star_{\mathrm{cut}}G)_m$.

The coherence in the
lemma supplies the coherence for composite surjections, so the
product is a word coefficient system.

For three factors, either parenthesization is a direct sum over
$0\leq p\leq q\leq n$. Each summand starts with the external
product on the three consecutive blocks. Restrict it to the locus
where values in different blocks are distinct, then apply open
direct image. Successive cut conditions impose precisely this
common open. Base change and composition of open direct images
identify the two expressions. Empty blocks impose no condition.

The same description holds for any number of factors. All
associativity diagrams, including the pentagon, restrict to the
identity of this common expression. The functorial coherences
therefore make them commute. With $K$, only the empty-block
summand survives on the appropriate side, giving the unit maps
and their compatibility.
\end{proof}

No permutation of the blocks has been introduced. In particular,
this construction supplies an associative product without a chosen
symmetry. The empty unit is distinct from a local unit field.
A local vacuum in a boundary algebra is extra data on nonempty
configurations, whereas $K_1=0$.

\begin{example}\label{ex:vol2-outer-cut}
At three positions,
\[
 U_{1,3}=X^3\setminus(\Delta_{12}\cup\Delta_{13}),\qquad
 U_{2,3}=X^3\setminus(\Delta_{13}\cup\Delta_{23}),
\]
where $\Delta_{ij}$ is the locus $x_i=x_j$. Their union has
complement exactly $\Delta_{13}$. If $x_1\ne x_3$ and the first
cut fails, then $x_1=x_2$, so the second cut succeeds.
Conversely, $x_1=x_3$ defeats both cuts. Thus every nonadjacent
equality crossing a cut is required already in the first
three-position calculation.

On an affine coordinate chart,
$u=z_1-z_2$ and $v=z_2-z_3$ give the outer equation $u+v=0$.
The stable and finite complexes of
Section~\ref{sec:vol2-supported-tensors} compute different
coefficient constructions on this same arrangement.
\end{example}

Call $F$ \emph{reduced} if $F_0=0$.
For reduced objects $F^{(1)},\ldots,F^{(r)}$, every nonzero
summand of their iterated cut product at $n$ has $r$ nonempty
consecutive blocks. Hence
\begin{equation}\label{eq:vol2-positive-word-length}
 (F^{(1)}\star_{\mathrm{cut}}\cdots
          \star_{\mathrm{cut}}F^{(r)})_n=0\qquad(r>n).
\end{equation}
This bound concerns word length. It makes tensor-length constructions
finite at a fixed configuration size. It does not bound the order
of a pole or the normal derivatives in a supported coefficient.

\section{Comparing attachments}
\label{sec:vol2-coefficient-attachments}

The localization sequence \eqref{eq:vol2-localization-sequence}
ties regular, meromorphic and supported coefficients together.
A comparison must respect this attachment in addition to the
separate objects. Let $r:U\to Z$ and $r':U'\to Z'$ be cochain
maps of complexes of left $D$-modules, and take $D$-linear Hom
complexes. For derived maps among $A,L,E$, use the bounded free
resolutions of Proposition~\ref{prop:vol2-spencer-maps}.

Give each cone the two-step filtration
$0\subset Z\subset\operatorname{Cone}(r)$, with the convention
of \eqref{eq:vol2-coefficient-cone}. A degree-zero map preserving
these filtrations has the form
\begin{equation}\label{eq:vol2-attachment-map}
 (\zeta,\xi)\longmapsto(g\zeta+H\xi,f\xi),\qquad
 |f|=|g|=0,\quad |H|=-1.
\end{equation}
Here $f:U\to U'$, $g:Z\to Z'$ and $H:U\to Z'$.

\begin{lemma}\label{lem:vol2-attachment-homotopy}
The map \eqref{eq:vol2-attachment-map} is closed if and only if
$f,g$ are cochain maps and
\begin{equation}\label{eq:vol2-attachment-homotopy}
 d_{\Hom}H=gr-r'f.
\end{equation}
The complex of filtration-preserving cone maps is isomorphic,
with the shifted-coordinate signs, to the homotopy fibre of
\[
 \Hom(U,U')\oplus\Hom(Z,Z')\longrightarrow\Hom(U,Z'),
 \qquad(f,g)\longmapsto gr-r'f.
\]
For closed degree-zero maps, composition is
\begin{equation}\label{eq:vol2-attachment-composition}
 (f_2,g_2,H_2)(f_1,g_1,H_1)
 =(f_2f_1,g_2g_1,g_2H_1+H_2f_1).
\end{equation}
\end{lemma}

\begin{proof}
Insert \eqref{eq:vol2-attachment-map} into the cone differential
$(\zeta,\xi)\mapsto(d_Z\zeta+r\xi,-d_U\xi)$. The two diagonal entries in the
chain-map equation say $df=dg=0$. The off-diagonal entry says
$d_{Z'}H+Hd_U=gr-r'f$, which is
\eqref{eq:vol2-attachment-homotopy}.

In general degree $j$, put a factor $(-1)^j$ on the $U[1]\to U'[1]$
entry and a factor $(-1)^j$ on the $U[1]\to Z'$ entry.
In these coordinates the differential is
\[
 (f,g,H)\longmapsto
 (df,dg,gr-r'f-dH).
\]
This is the homotopy-fibre differential, since
$|H|=j-1$. Finally, multiplication of the degree-zero upper triangular
block matrices gives \eqref{eq:vol2-attachment-composition} and
its associativity.
\end{proof}

Equations~\eqref{eq:vol2-connection-hom} and
\eqref{eq:vol2-attachment-homotopy} retain two kinds of compatibility:
with differentiation and with the residue attachment. The cut product
adds the geometry of ordered insertion. Together they specify
coefficient objects and comparison complexes for local operations.

An action on them must still give operations in every arity and
homotopies compatible with substitution and collision restriction.
The construction of a field theory must, independently, supply its
fields, differential, interaction and quantum observables.
It must also give a map from those observables to the geometric actions.
The next calculation
begins this theory-dependent part with finite Hamiltonian jets.
